\section{Introducción}
La realimentación completa de estados permite modificar los polos del sistema,
pero requiere conocer cada variable del vector de estado. Cuando sólo es posible
medir la salida de la planta, el diseño digital puede incluir un observador
capaz de reconstruir estados no medidos a partir del modelo, el mando y la
señal disponible \cite{samanta2024,fadali2020,astrom2020}.

La guía de esta práctica exige diseñar e implementar dos estructuras: un
estimador de predicción y un estimador de estado actual. Ambas se aplicarán a
una planta analógica de segundo orden. A diferencia del laboratorio previo,
la realimentación debe efectuarse \emph{únicamente sobre los estados estimados};
la medición física usada por el algoritmo será la salida del sistema.

La experiencia acumulada en el Laboratorio 5 proporciona una infraestructura
experimental de PSoC 5LP, MATLAB y adquisición con el Rigol DS1202Z-E. Sin
embargo, la calibración de esa infraestructura y el modelo de la nueva práctica
requieren verificación antes de cerrar el lazo. Esta distinción permite preservar
trazabilidad sin trasladar valores antiguos como resultados actuales.
